\documentclass[12pt,a4paper]{article}

% ── Paquetes ──────────────────────────────────────────────────────────────────
\usepackage[utf8]{inputenc}
\usepackage[T1]{fontenc}
\usepackage[spanish]{babel}
\usepackage{amsmath,amssymb,amsthm}
\usepackage{geometry}
\usepackage{graphicx}
\usepackage{booktabs}
\usepackage{hyperref}
\usepackage{listings}
\usepackage{xcolor}
\usepackage{parskip}
\usepackage{enumitem}
\usepackage{fancyhdr}
\usepackage{titlesec}

\geometry{margin=2.5cm, headheight=15pt}

% ── Estilo de código ──────────────────────────────────────────────────────────
\definecolor{codebg}{RGB}{245,245,245}
\definecolor{codeframe}{RGB}{200,200,200}
\definecolor{keywordcolor}{RGB}{0,100,180}
\definecolor{commentcolor}{RGB}{100,120,100}
\definecolor{stringcolor}{RGB}{160,40,40}

\lstset{
  language=Python,
  backgroundcolor=\color{codebg},
  frame=single,
  rulecolor=\color{codeframe},
  basicstyle=\ttfamily\small,
  keywordstyle=\color{keywordcolor}\bfseries,
  commentstyle=\color{commentcolor}\itshape,
  stringstyle=\color{stringcolor},
  numbers=left,
  numberstyle=\tiny\color{gray},
  numbersep=8pt,
  breaklines=true,
  showstringspaces=false,
  tabsize=4,
}

% ── Encabezado ────────────────────────────────────────────────────────────────
\pagestyle{fancy}
\fancyhf{}
\rhead{Modelado y Simulación}
\lhead{Filtro de Kalman — Señales de Voz}
\cfoot{\thepage}

% ── Título ────────────────────────────────────────────────────────────────────
\title{
  \vspace{-1cm}
  {\Large \textbf{Caso de Estudio}} \\[0.4em]
  {\LARGE Reconstrucción de Señales de Voz en Entornos Ruidosos} \\[0.3em]
  {\large mediante el Filtro de Kalman} \\[0.5em]
  \hrule
}
\author{Juan Bautista Espino \\[0.2em]
  \small Modelado y Simulación --- 2026, 1\textsuperscript{er} cuatrimestre}
\date{\today}

% ──────────────────────────────────────────────────────────────────────────────
\begin{document}
\maketitle
\tableofcontents
\newpage

% ══════════════════════════════════════════════════════════════════════════════
\section{Planteamiento del Problema}
% ══════════════════════════════════════════════════════════════════════════════

En aplicaciones reales de comunicación y procesamiento del habla —clases virtuales,
llamadas telefónicas, sistemas de reconocimiento automático de voz (ASR)—, la señal
captada por el micrófono no es la voz pura del locutor sino una versión degradada por:

\begin{itemize}[leftmargin=1.5cm]
  \item \textbf{Ruido ambiental}: tráfico, viento, conversaciones de fondo.
  \item \textbf{Ruido electrónico}: cuantificación del ADC, interferencias del circuito.
  \item \textbf{Pérdidas de paquetes}: en comunicación VoIP o streaming.
\end{itemize}

Estas perturbaciones reducen la inteligibilidad, dificultan el análisis automático y
degradan la experiencia perceptual del oyente.

\subsection*{Pregunta central}

\begin{center}
\textit{¿Cómo reconstruir una señal de voz más cercana a la original
a partir de mediciones ruidosas, usando herramientas de estimación estadística?}
\end{center}

% ══════════════════════════════════════════════════════════════════════════════
\section{Objetivo e Hipótesis}
% ══════════════════════════════════════════════════════════════════════════════

\subsection{Objetivo}

Evaluar cuantitativamente si el filtro de Kalman logra:
\begin{enumerate}
  \item Reducir el ruido aditivo en señales de voz sintéticas.
  \item Aproximar mejor la señal ``real'' que métodos de interpolación clásicos.
  \item Ofrecer una ventaja medible en las métricas MSE y SNR.
\end{enumerate}

\subsection{Hipótesis}

\textit{El filtro de Kalman proporciona una reconstrucción más fiel de la señal de
voz que la interpolación polinómica y el spline cúbico, gracias a su capacidad de
modelar explícitamente la incertidumbre del proceso y de la medición, e incorporar
la dinámica temporal de la señal.}

% ══════════════════════════════════════════════════════════════════════════════
\section{Marco Teórico}
% ══════════════════════════════════════════════════════════════════════════════

\subsection{Modelo de señal observada}

Se modela la situación como un sistema de estado discreto:

\begin{align}
  x_k &= x_{k-1} + w_k, \qquad w_k \sim \mathcal{N}(0,\,Q) \label{eq:estado}\\
  z_k &= x_k + v_k, \qquad\;\; v_k \sim \mathcal{N}(0,\,R) \label{eq:medicion}
\end{align}

donde $x_k$ es la amplitud real (desconocida) de la voz en el instante $k$,
$z_k$ es la muestra ruidosa observada, $w_k$ es el ruido de proceso (variabilidad
intrínseca de la señal) y $v_k$ es el ruido de medición (ruido ambiental/electrónico).
Las varianzas $Q$ y $R$ cuantifican la incertidumbre en cada fuente.

\subsection{Filtro de Kalman escalar}

El filtro de Kalman resuelve el problema de \textbf{estimación MMSE}
(Minimum Mean Square Error) para el modelo lineal gaussiano de las
ecuaciones \eqref{eq:estado}--\eqref{eq:medicion}.
Opera en dos pasos recursivos:

\subsubsection*{Paso 1 — Predicción}

Se proyecta el estado y la covarianza del error al instante siguiente sin
utilizar aún la medición $z_k$:

\begin{align}
  \hat{x}_{k \mid k-1} &= \hat{x}_{k-1 \mid k-1} \label{eq:pred_x}\\
  P_{k \mid k-1}        &= P_{k-1 \mid k-1} + Q   \label{eq:pred_P}
\end{align}

\subsubsection*{Paso 2 — Corrección}

Se incorpora la medición $z_k$ ponderando por la \textbf{ganancia de Kalman} $K_k$:

\begin{align}
  K_k &= \frac{P_{k \mid k-1}}{P_{k \mid k-1} + R} \label{eq:ganancia}\\
  \hat{x}_{k \mid k} &= \hat{x}_{k \mid k-1} + K_k \bigl(z_k - \hat{x}_{k \mid k-1}\bigr) \label{eq:correc_x}\\
  P_{k \mid k} &= (1 - K_k)\,P_{k \mid k-1} \label{eq:correc_P}
\end{align}

\paragraph{Interpretación de la ganancia.}
Si $R \ll P_{k|k-1}$ (alta confianza en la medición), entonces $K_k \to 1$
y el filtro sigue de cerca a $z_k$.
Si $Q \ll R$ (proceso muy suave y medición muy ruidosa), entonces $K_k \to 0$
y el filtro ignora la medición y confía en su predicción.

\subsection{Métodos clásicos de referencia}

\subsubsection*{Interpolación polinómica}

Dado el conjunto de puntos $\{(t_k, z_k)\}_{k=0}^{N-1}$, se busca el polinomio
de grado $n$ que minimiza el error cuadrático total (ajuste por mínimos cuadrados):

\begin{equation}
  P(t) = \sum_{j=0}^{n} c_j\, t^j, \qquad
  \mathbf{c} = \arg\min_{\mathbf{c}} \sum_{k=0}^{N-1}\bigl(z_k - P(t_k)\bigr)^2
\end{equation}

Para grados altos y datos equiespaciados aparecen las \textbf{oscilaciones de Runge}
en los extremos del intervalo, que amplifican el error en lugar de reducirlo.

\subsubsection*{Spline cúbico natural}

Se divide el dominio en subintervalos mediante nodos $\{t_{i_0}, t_{i_1}, \ldots\}$
obtenidos por submuestreo (cada $s$ muestras). En cada subintervalo se ajusta un
polinomio cúbico, imponiendo continuidad de primera y segunda derivada:

\begin{equation}
  S_j(t) = a_j + b_j(t - t_j) + c_j(t-t_j)^2 + d_j(t-t_j)^3,
  \qquad t \in [t_j, t_{j+1}]
\end{equation}

con condiciones de frontera naturales $S''(t_0) = S''(t_{N-1}) = 0$.

\subsection{Métricas de evaluación}

\begin{align}
  \mathrm{MSE} &= \frac{1}{N}\sum_{k=1}^{N}\bigl(x_k - \hat{x}_k\bigr)^2
  \label{eq:mse}\\[4pt]
  \mathrm{SNR} &= 10\log_{10}\!\left(\frac{\displaystyle\sum_{k}x_k^2}
  {\displaystyle\sum_{k}(x_k-\hat{x}_k)^2}\right) \quad [\mathrm{dB}]
  \label{eq:snr}
\end{align}

Un MSE menor y un SNR mayor indican mejor reconstrucción.

% ══════════════════════════════════════════════════════════════════════════════
\section{Metodología}
% ══════════════════════════════════════════════════════════════════════════════

\subsection{Generación de la señal de voz sintética}

Se genera una señal tipo voz mediante una suma de $H = 7$ armónicos con amplitudes
decrecientes (modelo de fuente glótica simplificado) y una envolvente de amplitud lenta:

\begin{equation}
  x(t) = A(t) \cdot \sum_{k=1}^{H} a_k \sin\!\bigl(2\pi k f_0 t + \varphi_k\bigr),
  \qquad A(t) = \tfrac{1}{2}\!\left(1 + \sin\!\bigl(2\pi \cdot 1.5\cdot t + \tfrac{\pi}{4}\bigr)\right)
\end{equation}

con $f_0 = 150$ Hz (frecuencia fundamental), $a_k = \{1,\,0.6,\,0.35,\,0.2,\,0.12,\,0.07,\,0.04\}$
y fases $\varphi_k$ aleatorias uniformes en $[0, 2\pi)$.
La señal se normaliza al rango $[-1, 1]$ y se muestrea a $f_s = 8000$ Hz.

\subsection{Adición de ruido blanco gaussiano}

Para simular un entorno con relación señal/ruido de entrada $\mathrm{SNR_{in}}$ dB,
se añade ruido blanco $v \sim \mathcal{N}(0, \sigma_v^2)$ donde:

\begin{equation}
  \sigma_v^2 = \frac{P_x}{10^{\,\mathrm{SNR_{in}}/10}}, \qquad P_x = \frac{1}{N}\sum_k x_k^2
\end{equation}

\subsection{Pipeline de comparación}

\begin{enumerate}
  \item Generar $x$ (señal limpia) y $z = x + v$ (señal ruidosa).
  \item Aplicar \textbf{interpolación polinómica} (grado $n = 8$) sobre $z$.
  \item Aplicar \textbf{spline cúbico} (nodos cada $s = 5$ muestras) sobre $z$.
  \item Aplicar \textbf{filtro de Kalman} ($Q = 10^{-2}$, $R = 10^{-2}$) sobre $z$.
  \item Calcular MSE y SNR de salida de cada método comparando con $x$.
\end{enumerate}

\subsection{Implementación en Python}

El módulo \texttt{kalman\_voz.py} implementa todos los métodos.
A continuación se muestra el núcleo del filtro de Kalman:

\begin{lstlisting}[caption={Filtro de Kalman escalar — \texttt{kalman\_voz.py}}]
def filtro_kalman(z, Q=1e-2, R=1e-2, x0=0.0, P0=1.0):
    n = len(z)
    x_hat = np.zeros(n)
    x_pred = x0
    P = P0

    for k in range(n):
        # -- Prediccion --
        P_pred = P + Q

        # -- Correccion --
        K = P_pred / (P_pred + R)          # Ganancia de Kalman
        x_hat[k] = x_pred + K * (z[k] - x_pred)
        P = (1 - K) * P_pred

        x_pred = x_hat[k]

    return x_hat
\end{lstlisting}

La generación de la señal sintética:

\begin{lstlisting}[caption={Generación de voz sintética — \texttt{kalman\_voz.py}}]
def generar_senal_voz(n_muestras=400, fs=8000.0, freq_fundamental=150.0, seed=42):
    rng = np.random.default_rng(seed)
    t = np.linspace(0, n_muestras / fs, n_muestras, endpoint=False)

    x = np.zeros(n_muestras)
    for k, amp in enumerate([1.0, 0.6, 0.35, 0.2, 0.12, 0.07, 0.04], start=1):
        phase = rng.uniform(0, 2 * np.pi)
        x += amp * np.sin(2 * np.pi * k * freq_fundamental * t + phase)

    env = 0.5 * (1 + np.sin(2 * np.pi * 1.5 * t + np.pi / 4))
    x *= env
    x /= np.max(np.abs(x) + 1e-12)
    return t, x
\end{lstlisting}

% ══════════════════════════════════════════════════════════════════════════════
\section{Resultados}
% ══════════════════════════════════════════════════════════════════════════════

\subsection{Configuración del experimento}

\begin{center}
\begin{tabular}{ll}
\toprule
\textbf{Parámetro} & \textbf{Valor} \\
\midrule
Muestras ($N$)             & 400 \\
Frecuencia de muestreo ($f_s$) & 8\,000 Hz \\
Frecuencia fundamental ($f_0$) & 150 Hz \\
SNR de entrada             & 10 dB \\
Grado del polinomio        & 8 \\
Factor de submuestreo (spline) & 5 \\
$Q$ (Kalman)               & $10^{-2}$ \\
$R$ (Kalman)               & $10^{-2}$ \\
\bottomrule
\end{tabular}
\end{center}

\subsection{Métricas obtenidas}

\begin{center}
\begin{tabular}{lcc}
\toprule
\textbf{Método} & \textbf{MSE} & \textbf{SNR de salida (dB)} \\
\midrule
Señal ruidosa (sin filtrar) & 0.016644 & 10.43 \\
Interpolación polinómica    & 0.169189 &  0.36 \\
Spline cúbico               & 0.017644 & 10.18 \\
\textbf{Filtro de Kalman}   & \textbf{0.013317} & \textbf{11.79} \\
\bottomrule
\end{tabular}
\end{center}

\subsection{Análisis de resultados}

\paragraph{Interpolación polinómica.}
El polinomio de grado 8 presenta el peor desempeño: MSE $= 0.169$, SNR $= 0.36$ dB.
Al ajustar globalmente sobre los datos ruidosos, el método introduce las oscilaciones
de Runge en los extremos del intervalo, \emph{amplificando} el error en lugar de
reducirlo. Este comportamiento confirma que la interpolación global es inadecuada
para señales estocásticas con ruido.

\paragraph{Spline cúbico.}
Al utilizar nodos submuestreados (1 de cada 5 muestras), el spline evita interpolar
exactamente el ruido, obteniendo MSE $= 0.018$ y SNR $= 10.18$ dB, levemente inferior
a la señal ruidosa sin procesar. Aunque el suavizado local es conceptualmente correcto,
el método no tiene mecanismo para estimar cuánto del residuo es ruido real versus
variación legítima de la señal.

\paragraph{Filtro de Kalman.}
Con $Q = R = 10^{-2}$, el filtro obtiene \textbf{MSE $= 0.013$} y \textbf{SNR $= 11.79$ dB},
superando a todos los métodos. La mejora de $+1.36$ dB sobre la señal ruidosa y de
$+1.61$ dB sobre el spline es consecuencia directa de la estimación probabilística:
en cada paso, el filtro pondera la predicción del modelo con la observación actual
según la relación $K_k = P_{k|k-1}/(P_{k|k-1}+R)$, asignando mayor peso a la fuente
más confiable.

% ══════════════════════════════════════════════════════════════════════════════
\section{Discusión}
% ══════════════════════════════════════════════════════════════════════════════

\subsection{¿Por qué el Filtro de Kalman supera a la interpolación?}

Los métodos de interpolación son esencialmente \textbf{deterministas}: buscan una
curva que pase ``cerca'' de los datos observados, sin distinguir si la variabilidad
es parte de la señal o ruido. El filtro de Kalman, en cambio, opera en un marco
\textbf{probabilístico}: modela explícitamente dos fuentes de incertidumbre ($Q$ y $R$)
y computa en cada instante la estimación que minimiza el error cuadrático esperado.

\subsection{Rol de los parámetros Q y R}

La razón $Q/R$ controla el comportamiento del filtro:

\begin{itemize}
  \item $Q \ll R$: el filtro confía más en su predicción que en la medición.
    Produce una salida muy suavizada pero potencialmente retrasada o sesgada.
  \item $Q \gg R$: el filtro sigue de cerca cada muestra nueva.
    En el límite ($R \to 0$), la salida coincide con la señal ruidosa.
  \item $Q \approx R$: equilibrio óptimo para señales moderadamente no estacionarias
    con SNR de entrada $\sim 10$ dB.
\end{itemize}

\subsection{Limitaciones del modelo de paseo aleatorio}

El modelo de estado \eqref{eq:estado} asume que la variación de $x_k$ entre muestras
consecutivas es ruido blanco. Para señales de voz reales, esto es una aproximación:
la voz tiene correlación de largo plazo (prosodia, vocales sostenidas) que un modelo
autoregresivo (AR) o un Kalman de orden superior capturaría mejor.
Sin embargo, para el propósito de este caso de estudio, el modelo escalar demuestra
con claridad las ventajas del enfoque probabilístico.

% ══════════════════════════════════════════════════════════════════════════════
\section{Conclusiones}
% ══════════════════════════════════════════════════════════════════════════════

\begin{enumerate}
  \item El filtro de Kalman escalar, con parámetros $Q = R = 10^{-2}$, logró
        \textbf{SNR $= 11.79$ dB}, superando la señal ruidosa ($10.43$ dB),
        el spline cúbico ($10.18$ dB) y la interpolación polinómica ($0.36$ dB).

  \item La interpolación polinómica global es \textbf{contraindicada} para señales
        estocásticas con ruido: las oscilaciones de Runge empeoran la reconstrucción
        respecto a no aplicar ningún filtrado.

  \item El spline cúbico ofrece una mejora marginal frente a la señal sin procesar,
        pero no puede explotar la estructura probabilística del problema.

  \item La hipótesis planteada queda \textbf{confirmada}: el filtro de Kalman supera
        a los métodos clásicos gracias a su modelado explícito del ruido de proceso
        y de medición, y a su actualización recursiva que incorpora la dinámica temporal.

  \item Las líneas de trabajo futuro incluyen: modelos AR de orden superior,
        filtro de Kalman extendido para señales no lineales, e integración con
        pipelines de reconocimiento automático de voz (ASR).
\end{enumerate}

\vspace{1em}
\begin{center}
\textit{``Mientras los métodos de interpolación intentan describir los datos observados,}\\
\textit{el filtro de Kalman busca inferir la realidad que los genera.''}
\end{center}

% ══════════════════════════════════════════════════════════════════════════════
\section*{Referencias}
% ══════════════════════════════════════════════════════════════════════════════

\begin{enumerate}
  \item Kalman, R.~E. (1960). \textit{A new approach to linear filtering and prediction
        problems}. Journal of Basic Engineering, 82(1), 35--45.
  \item Haykin, S. (2001). \textit{Kalman Filtering and Neural Networks}.
        Wiley-Interscience.
  \item Proakis, J.~G., \& Manolakis, D.~G. (2006). \textit{Digital Signal Processing:
        Principles, Algorithms, and Applications} (4th ed.). Prentice Hall.
  \item Virtanen, P. et al. (2020). \textit{SciPy 1.0: Fundamental Algorithms for
        Scientific Computing in Python}. Nature Methods, 17, 261--272.
\end{enumerate}

\end{document}
